% ================================================================== Chapitre 0
\chapter*{Programme, plan et notations}
\addcontentsline{toc}{chapter}{Programme, plan et notations}
\markboth{Programme, plan et notations}{}

\section*{Ce que le programme exige}

Ce chapitre regroupe trois blocs du programme de PT. Chaque capacité ci-dessous est traitée dans le cours (le numéro du chapitre est indiqué entre crochets) et au moins un exercice corrigé la mobilise.

\begin{resultat}{Capacités exigibles -- résumé}
\textbf{Premier principe appliqué à la transformation chimique}
\begin{itemize}
  \item calculer une enthalpie standard de réaction avec des tables ou la loi de Hess ; relier son signe au caractère endo- ou exothermique \hfill[3]
  \item prévoir le sens et estimer le transfert thermique en réacteur monobare isotherme \hfill[4]
  \item évaluer la température atteinte en réacteur adiabatique monobare ; simuler $T(t)$ en Python \hfill[4]
\end{itemize}
\textbf{Deuxième principe appliqué à la transformation chimique}
\begin{itemize}
  \item prévoir le signe de l'entropie standard de réaction \hfill[3]
  \item relier création d'entropie et enthalpie libre de réaction à $T$ et $P$ fixées \hfill[5]
  \item prévoir le sens d'évolution avec $\Delta_r G$ (ou $Q_r$ et $K^\circ$) \hfill[5]
  \item calculer les grandeurs standard de réaction (tables, Hess, combinaisons linéaires) \hfill[3]
  \item citer et exploiter la relation de Van~'t~Hoff ; calculer $K^\circ(T)$ à toute température \hfill[6]
  \item déterminer l'état final : équilibre chimique ou transformation totale \hfill[5]
  \item identifier les paramètres d'optimisation d'une synthèse \hfill[6]
\end{itemize}
\textbf{Thermodynamique de l'oxydoréduction et piles}
\begin{itemize}
  \item relier $\Delta_rG$ aux potentiels des couples et $\Delta_rG^\circ$ aux potentiels standard \hfill[7]
  \item calculer un potentiel standard inconnu à partir de données thermodynamiques \hfill[7]
  \item établir l'inégalité entre variation d'enthalpie libre et travail électrique \hfill[8]
  \item relier la tension à vide d'une pile à l'enthalpie libre de réaction ; calculer sa capacité \hfill[8]
\end{itemize}
\end{resultat}

\noindent Trois précisions du programme guident tout le cours :
\begin{itemize}
  \item l'expression $\mu_i=\mu_i^\circ+RT\ln a_i$ du potentiel chimique est \emph{admise} (sa démonstration est hors programme) ;
  \item on suppose $\rHo$ et $\rSo$ \emph{indépendantes de la température} en dehors des changements d'état (approximation d'Ellingham) ;
  \item la notion d'\emph{affinité chimique} n'est pas utilisée : le sens d'évolution se lit sur le signe de $\Delta_rG$.
\end{itemize}

\section*{Notations utilisées}

\renewcommand{\arraystretch}{1.25}
\begin{xltabular}{\linewidth}{>{\centering\arraybackslash}p{2.2cm} X >{\centering\arraybackslash}p{3.2cm}}
\toprule
\textbf{Symbole} & \textbf{Signification} & \textbf{Unité} \\ \midrule
$\nu_i$ & coefficient stœchiométrique \emph{algébrique} ($>0$ produit, $<0$ réactif) & sans \\
$\xi$ & avancement de la réaction & mol \\
$a_i$ & activité du constituant $\mathrm{A}_i$ & sans \\
$\mu_i$, $\mu_i^\circ$ & potentiel chimique, potentiel chimique standard & $\mathrm{J\cdot mol^{-1}}$ \\
$G$, $H$, $S$ & enthalpie libre, enthalpie, entropie du système & J, J, $\mathrm{J\cdot K^{-1}}$ \\
$\Delta_r X$ & grandeur de réaction ($X = G, H, S$) & $\mathrm{J\cdot mol^{-1}}$ (ou $\mathrm{J\cdot K^{-1}\cdot mol^{-1}}$) \\
$\Delta_r X^\circ$ & grandeur standard de réaction & idem \\
$\Delta_f X^\circ$ & grandeur standard de formation & idem \\
$Q_r$ & quotient de réaction & sans \\
$K^\circ$ & constante thermodynamique d'équilibre (notée aussi $K$) & sans \\
$Q$ & transfert thermique reçu par le système & J \\
$\gamma$ & nombre d'électrons échangés dans la réaction redox & sans \\
$E$, $E^\circ$ & potentiel d'électrode (de Nernst), potentiel standard & V \\
$e$, $e^\circ$ & tension à vide d'une pile ($e=E_+-E_-$), tension à vide standard & V \\
$q$, $q_{\max}$ & charge débitée, capacité d'une pile & C (ou A$\cdot$h) \\
\bottomrule
\end{xltabular}

\begin{ordres}[constantes utiles]
\begin{itemize}
  \item constante des gaz parfaits : $R = \SI{8,314}{J.K^{-1}.mol^{-1}}$ ; constante de Faraday : $F = \mathcal{N}_A e = \SI{96485}{C.mol^{-1}} \approx \SI{96500}{C.mol^{-1}}$ ;
  \item pression standard : $P^\circ = \SI{1}{bar} = \SI{e5}{Pa}$ ; concentration standard : $c^\circ = \SI{1}{mol.L^{-1}}$ ;
  \item à $T = \SI{298}{K}$ : $RT = \SI{2,48}{kJ.mol^{-1}}$ et
  \[ \frac{RT\ln 10}{F} = \SI{0,0592}{V} \approx \SI{0,06}{V} \qquad\text{(c'est le « 0,06 » des formules de Nernst).} \]
\end{itemize}
\end{ordres}
